\chapter{Einleitung}
\label{cha:Einleitung}

Für funktionierende zusammenhängende Produktionsanlagen ist der Austausch vieler Daten zwischen einzelnen Maschinen essentiell. Gerade in komplexen Verpackungslinien der Firma MULTIVAC werden ständig Daten über Zustände, Prozesse und verschiedene Ereignisse ausgetauscht, wodurch ein störungsfreier Produktionsablauf sichergestellt wird.
Bei MULTIVAC wird für diese Kommunikation das Nachrichtenprotokoll Message Queuing Telemetry Transport (MQTT) eingesetzt.

Mit steigenden Anforderungen an Verpackungsmaschinen und zunehmender Anlagenkomplexität wächst auch das Nachrichtenaufkommen stark an. Dadurch wird die Analyse und Fehlerdiagnose der Kommunikationsabläufe jedoch zunehmend schwieriger. Bestehende Werkzeuge ermöglichen zwar die Anzeige einzelner Nachrichten, bieten jedoch nur begrenzte Unterstützung bei der systematischen Fehlersuche und Diagnose. Dies führt zu einem hohen Analyseaufwand und erfordert umfangreiches Expertenwissen für die speziell in MULTIVAC entwickelte MQTT-Umsetzung.

Ziel dieser Projektarbeit ist daher die Entwicklung einer Überwachungs- und Diagnose-Applikation zur Analyse der MQTT-basierten Maschinenkommunikation in Verpackungslinien. Die Anwendung soll dabei Kommunikationsabläufe transparent darstellen, Nachrichten dauerhaft speichern sowie eine effiziente Analyse großer Datenmengen ermöglichen. Dadurch soll der Analyseaufwand reduziert werden. Darüber hinaus sollen Kommunikations- und Prozessfehler automatisch erkannt sowie Funktionen zur Teilnehmerüberwachung und zur manuellen Nachrichtensimulation bereitgestellt werden. Dadurch soll der Arbeitsaufwand für die Fehlersuche bei der Inbetriebnahme verringert werden.